% Generated draft: requires mathematical and source review.
\documentclass[11pt]{article}
\usepackage[margin=25mm]{geometry}
\usepackage{fontspec}
\setmainfont{Latin Modern Roman}
\usepackage{amsmath,amssymb,mathtools,unicode-math}
\setmathfont{Latin Modern Math}
\usepackage{enumitem,needspace,xurl,hyperref}
\setlength{\parindent}{0pt}
\setlength{\parskip}{5pt}
\setlength{\emergencystretch}{3em}
\newcommand{\N}{\mathbb{N}}
\newcommand{\Z}{\mathbb{Z}}
\newcommand{\Q}{\mathbb{Q}}
\newcommand{\R}{\mathbb{R}}
\newcommand{\C}{\mathbb{C}}
\newcommand{\F}{\mathbb{F}}
\newcommand{\A}{\mathbb{A}}
\newcommand{\PP}{\mathbb{P}}
\newcommand{\E}{\mathbb{E}}
\newcommand{\eps}{\varepsilon}
\newcommand{\dd}{\,\mathrm{d}}
\begin{document}
\section*{UnsolvedMath ID: 30002011}
\textbf{Source problem code:} catalogue entry
\subsection{Definitions and mathematical statement}
Let \(\Delta_h\) denote the monotoneized empirical-Bayes estimator referred to in the supplied source record, with corruption parameter \(h\). The record does not provide a formula for \(\Delta_h\), a definition of the corruption operation or monotoneization, or the admissible values of \(h\). The term \emph{inbred cross-validation} denotes the proposed selection method named in the source; its construction is not supplied.

The source asks:
\begin{enumerate}
\item How should the corruption parameter \(h\) be chosen for \(\Delta_h\)?
\item Does choosing \(h\) by the proposed inbred cross-validation method yield a rate-sharp procedure?
\end{enumerate}

These questions retain unresolved qualifications. The observation model, parameter space, loss, risk or regret criterion, comparison class, asymptotic regime, and meaning of rate sharpness are unspecified. Consequently, the supplied record does not determine a fully quantified mathematical assertion or a precise rate bound to be established. Its background mentions sharp regret rates for related Poisson empirical-Bayes procedures, but does not supply a Poisson model or a rate benchmark for this particular estimator and selector.
\subsection{Short English statement}
The problem asks how to tune the corruption parameter \(h\) in a monotoneized empirical-Bayes estimator \(\Delta_h\), and whether the proposed inbred cross-validation selector achieves a sharp rate. The estimator, selector, statistical model, and performance benchmark require specification before this becomes a precise mathematical conjecture.
\subsection{Sources}
\begin{itemize}
\item Local UnsolvedMath registry: supplied record with permanent ID \(30002011\), exact source problem code \texttt{catalogue entry}, titled \emph{Corruption-Parameter Choice in Empirical-Bayes Estimation}, in the Oberwolfach Reports Open Problems collection. The proposed year is 2012; no individual proposer is identified.
\item Supplied report citation: \emph{Frontiers in Nonparametric Statistics} (2012), Oberwolfach Reports, pp.~811--883, DOI \texttt{10.4171/owr/2012/14}. \url{https://doi.org/10.4171/owr/2012/14}. The supplied background describes Oberwolfach Reports as documenting mathematical workshops and their open-problem discussions.
\item Supplied literature reference: \emph{Optimal Poisson empirical Bayes estimation}. \url{https://arxiv.org/abs/2209.01328}.
\item Supplied literature assessment, checked on \texttt{2026-08-21}: the problem is labeled open; the proposed data-splitting cross-validation selector for \(h\) is described as unsettled. The assessment states that later work proves sharp regret rates for related Poisson empirical-Bayes procedures but does not resolve this specific selector. This is supplied dated curation, not an independently verified literature assessment; its time-sensitive status and cited claims require independent verification.
\item The supplied statement status is \texttt{corrected\_verified}. Its accompanying verification note states that the official report explicitly leaves the choice of \(h\) and the performance of inbred cross-validation open. This is a supplied source-verification claim, not an independent verification performed for this draft.
\end{itemize}
\end{document}
